% !TeX program = pdflatex
\documentclass[aspectratio=169,xcolor={dvipsnames,table}]{beamer}

\input{../../_en_preambulo_beamer}
% Comandos y macros estadísticas compartidas para las presentaciones Beamer en inglés (Metropolis)

% Conjuntos numéricos básicos
\newcommand{\R}{\mathbb{R}}
\newcommand{\N}{\mathbb{N}}
\newcommand{\Q}{\mathbb{Q}}
\newcommand{\Z}{\mathbb{Z}}
\newcommand{\set}[1]{\left\{ #1 \right\}}
\newcommand{\abs}[1]{\left|#1\right|}
\newcommand{\norm}[1]{\left\| #1 \right\|}

% Comandos estadísticos y probabilísticos del curso
\newcommand{\Var}{\operatorname{Var}}
\newcommand{\cov}{\operatorname{Cov}}
\newcommand{\corr}{\rho}
\newcommand{\comb}[2]{\begin{pmatrix} #1 \\ #2 \end{pmatrix}}
\newcommand{\s}{\sigma}
\newcommand{\particion}{\mathfrak{P}}
\newcommand{\card}[1]{\#\left( #1 \right)}
\newcommand{\seq}[3]{\set{#1}_{#2}^{#3}}
\newcommand{\E}{\mathbb{E}}
\newcommand{\Prob}{\mathbb{P}}

% Entornos pedagógicos y cajas en inglés académico natural (soportando tanto nombres en inglés como en español)
\newenvironment{discussion}{%
  \begin{alertblock}{Discussion Question}%
}{%
  \end{alertblock}%
}
\newenvironment{discusion}{%
  \begin{alertblock}{Discussion Question}%
}{%
  \end{alertblock}%
}

\newenvironment{reflection}{%
  \begin{alertblock}{Classroom Reflection}%
}{%
  \end{alertblock}%
}
\newenvironment{reflexion}{%
  \begin{alertblock}{Classroom Reflection}%
}{%
  \end{alertblock}%
}

\newenvironment{paradox}{%
  \begin{alertblock}{Exploratory Paradox}%
}{%
  \end{alertblock}%
}
\newenvironment{paradoja}{%
  \begin{alertblock}{Exploratory Paradox}%
}{%
  \end{alertblock}%
}

\newenvironment{motivation}{%
  \begin{exampleblock}{Motivation: Data Science in Practice}%
}{%
  \end{exampleblock}%
}
\newenvironment{motivacion}{%
  \begin{exampleblock}{Motivation: Data Science in Practice}%
}{%
  \end{exampleblock}%
}

\newenvironment{quickchallenge}{%
  \begin{block}{Quick Team Challenge}%
}{%
  \end{block}%
}
\newenvironment{retoclase}{%
  \begin{block}{Quick Team Challenge}%
}{%
  \end{block}%
}


\title{Chi-Squared Distribution and Sample Variance}
\subtitle{Section 05.03 --- Fisher's Theorem, Independence of Xbar and S\textsuperscript{2}, and CIs for Sigma\textsuperscript{2}}
\author[J. Castillo Colmenares]{Juliho Castillo Colmenares}
\institute[Tec de Monterrey]{Tecnologico de Monterrey}
\date{\vspace{-1.5cm}}

\begin{document}

% SLIDE 01: Title Page
\begin{frame}[plain]
	\titlepage
\end{frame}

% SLIDE 02: Roadmap
\begin{frame}{Roadmap --- Chapter 05: Sampling Distributions}
	\vspace{-0.15cm}
	\begin{block}{Curricular Structure of Unit 4}
		\footnotesize
		\begin{enumerate}\itemsep=0.03cm
			\item \textcolor{gray}{Section 05.01: Simple Random Sampling, Mean, and Unbiased Sample Variance}
			\item \textcolor{gray}{Section 05.02: Asymptotic Central Limit Theorem}
			\item \textbf{\textcolor{TecRojo}{Section 05.03: Chi-Squared Distribution and Sample Variance}} $\leftarrow$ \emph{Current focus}
			\item \textcolor{gray}{Section 05.04: Student's $t$-Distribution and Small Samples}
			\item \textcolor{gray}{Section 05.05: Fisher-Snedecor $F$-Distribution}
		\end{enumerate}
	\end{block}
	\vspace{-0.2cm}
	\begin{alertblock}{Learning Objective}
		\scriptsize
		Formalize the distribution of the sample variance $S^2$ for normal populations via Fisher's theorem, prove the surprising independence between $\bar{X}$ and $S^2$, and construct exact confidence intervals for $\sigma^2$.
	\end{alertblock}
\end{frame}

% SLIDE 03: Motivation
\begin{frame}{From Known $\sigma$ to Unknown $\sigma$}
	\footnotesize
	\begin{motivacion}
		The CLT (Section 05.02) approximates $\bar{X}$ with a Normal when $\sigma$ is \emph{known}. In practice we almost always must estimate $\sigma$ with $S$. For that we need to understand the exact distribution of $S^2$ --- and this is where the chi-squared comes in.
	\end{motivacion}
	\vspace{0.15cm}
	\pause
	\begin{itemize}\itemsep=0.1cm
		\item \textbf{Origin:} the chi-squared arises from summing squares of independent standard normals. \pause
		\item \textbf{Connection to 04.07:} we already saw that $\chi^2_\nu = \text{Gamma}(\nu/2, 2)$. \pause
		\item \textbf{Central question:} what is the exact distribution of $(n-1)S^2/\sigma^2$?
	\end{itemize}
\end{frame}

% SLIDE 03B: Definition and Properties
\begin{frame}{Definition and Properties of the Chi-Squared Distribution}
	\vspace{-0.15cm}
	\begin{columns}[T]
		\begin{column}{0.48\textwidth}
			\begin{block}{Definition and Density}
				\footnotesize
				$\chi^2_\nu = \sum_{i=1}^\nu Z_i^2$ for iid $Z_i\sim N(0,1)$. Density: $f(x)=\frac{1}{2^{\nu/2}\Gamma(\nu/2)}x^{\nu/2-1}e^{-x/2}$.
			\end{block}
		\end{column}
		\begin{column}{0.48\textwidth}
			\begin{alertblock}{Properties}
				\footnotesize
				$E(\chi^2_\nu)=\nu$, $\Var(\chi^2_\nu)=2\nu$; special case of Gamma$(\nu/2,2)$; reproductive: $\chi^2_{\nu_1}+\chi^2_{\nu_2}\sim\chi^2_{\nu_1+\nu_2}$.
			\end{alertblock}
		\end{column}
	\end{columns}
\end{frame}

% SLIDE 04: Fisher's Theorem
\begin{frame}{Fisher's Theorem}
	\vspace{-0.15cm}
	\begin{columns}[T]
		\begin{column}{0.48\textwidth}
			\begin{block}{Statement}
				\footnotesize
				For normal iid $X_1,\ldots,X_n$: (1) $\bar{X}$ and $S^2$ are independent; (2) $(n-1)S^2/\sigma^2 \sim \chi^2_{n-1}$.
			\end{block}
		\end{column}
		\begin{column}{0.48\textwidth}
			\begin{alertblock}{Why Surprising?}
				\footnotesize
				$S^2$ is built from $(X_i-\bar X)$, yet is independent of $\bar X$. An orthogonal transformation separates $\bar X$'s direction from the $n-1$ orthogonal directions capturing $(n-1)S^2$ (a case of Cochran's theorem).
			\end{alertblock}
		\end{column}
	\end{columns}
\end{frame}

% SLIDE 05: Confidence Interval for sigma^2
\begin{frame}{Exact Confidence Interval for $\sigma^2$}
	\vspace{-0.15cm}
	\begin{columns}[T]
		\begin{column}{0.48\textwidth}
			\begin{block}{Interval Formula}
				\footnotesize
				$\left[\dfrac{(n-1)S^2}{\chi^2_{n-1,1-\alpha/2}},\; \dfrac{(n-1)S^2}{\chi^2_{n-1,\alpha/2}}\right]$. Unlike Normal-based intervals, this one is \emph{asymmetric} around $S^2$.
			\end{block}
		\end{column}
		\begin{column}{0.48\textwidth}
			\begin{alertblock}{Hypothesis Testing for $\sigma^2$}
				\footnotesize
				The same statistic $(n-1)S^2/\sigma_0^2$ tests $H_0:\sigma^2=\sigma_0^2$ against $\chi^2_{n-1}$ quantiles.
			\end{alertblock}
		\end{column}
	\end{columns}
\end{frame}

% SLIDE 06: Applications
\begin{frame}{Applications in Inference and Quality Control}
	\vspace{-0.1cm}
	\begin{itemize}\itemsep=0.1cm
		\item \textbf{Quality control:} tests about manufacturing process variability against a target $\sigma^2$. \pause
		\item \textbf{Goodness-of-fit:} the $\chi^2$ test compares observed vs. expected frequencies. \pause
		\item \textbf{Preview 05.04:} the Student's $t$-distribution is built by combining a Normal with an independent chi-squared. \pause
		\item \textbf{Preview 05.05:} the $F$-distribution compares two variances via a ratio of two chi-squared variables.
	\end{itemize}
\end{frame}

% SLIDE 07: Python Lab Block 1
\begin{frame}[fragile]{Python Lab: Chi-Squared Properties}
	\vspace{-0.05cm}
	\scriptsize
	Mean/variance verification, the Gamma connection, and reproductivity (\texttt{05.03\_chi\_squared\_distribution.py}):
	{\lstset{basicstyle=\fontsize{4pt}{4.8pt}\selectfont\ttfamily,aboveskip=0.1em,belowskip=0.1em}
	\lstinputlisting[firstline=18, lastline=40]{../../code/05_distribuciones_muestreo/05.03_chi_squared_distribution.py}}
\end{frame}

% SLIDE 08: Python Lab Block 2
\begin{frame}[fragile]{Python Lab: Fisher's Theorem}
	\vspace{-0.05cm}
	\scriptsize
	Monte Carlo verification of the independence of $\bar X$ and $S^2$, and of $(n-1)S^2/\sigma^2\sim\chi^2_{n-1}$:
	{\lstset{basicstyle=\fontsize{4pt}{4.8pt}\selectfont\ttfamily,aboveskip=0.1em,belowskip=0.1em}
	\lstinputlisting[firstline=43, lastline=70]{../../code/05_distribuciones_muestreo/05.03_chi_squared_distribution.py}}
\end{frame}

% SLIDE 09: Python Lab Block 3
\begin{frame}[fragile]{Python Lab: Confidence Interval Coverage for $\sigma^2$}
	\vspace{-0.05cm}
	\scriptsize
	Empirical verification of the coverage of the chi-squared confidence interval for $\sigma^2$:
	{\lstset{basicstyle=\fontsize{4pt}{4.8pt}\selectfont\ttfamily,aboveskip=0.1em,belowskip=0.1em}
	\lstinputlisting[firstline=73, lastline=100]{../../code/05_distribuciones_muestreo/05.03_chi_squared_distribution.py}}
\end{frame}

% SLIDE 10: Python Lab Terminal Output
\begin{frame}[fragile]{Python Lab: Terminal Output}
	\vspace{-0.1cm}
	\begin{block}{}
		\fontsize{4pt}{5pt}\selectfont
		\begin{verbatim}
=== Block 1: Chi-Squared Properties ===
chi2_15: mean=15.00, variance=30.00 (matches Gamma(7.5, 2))
Additivity chi2_5+chi2_8 vs chi2_13: KS D=0.0013, p=0.908

=== Block 2: Fisher's Theorem ===
Correlation(Xbar, S^2) = 0.0038 (independence confirmed)
(n-1)*S^2/sigma^2 vs chi2_9: KS D=0.0018, p=0.517
Worked example: statistic=16.20, critical chi2_9,0.95=16.92 -> Reject H0? No

=== Block 3: Confidence Interval Coverage ===
chi2_19,0.025=8.91, chi2_19,0.975=32.85
Empirical coverage: 0.9487 (nominal: 0.9500)
95% CI for sigma^2 (n=20, S^2=45): [26.03, 96.00]
		\end{verbatim}
	\end{block}
\end{frame}



\end{document}
